\section{Checking the Lean examples}\label{appD:sec:lean}
Install Lean by following the official instructions~\cite{leaninstall}. As Section~\ref{ch15:sec:checker} explains, they set up the editor VS Code together with its Lean extension. Then create a file named \code{proofs.lean} with the same shape as in that section:
\par\noindent\begin{minipage}{\linewidth}
\begin{lstlisting}
import Std

namespace MathCourse

-- the definitions and theorems go here

end MathCourse
\end{lstlisting}
\end{minipage}\par
Lines beginning with \code{-{}-} are comments, which Lean ignores. Between the two lines that contain \code{MathCourse}, copy the Lean declarations of the book in the order in which they appear: first those of Chapter~\ref{ch:lean1}, then those of Chapter~\ref{ch:lean2}, and finally the Lean solutions for Weeks~15 and~16 in Appendix~\ref{app:solc}. A declaration is a block that begins with \code{def} or \code{theorem}. Some sections also display a proof state, part of a proof, or an error message; these show what Lean prints, and they are not meant to be copied. The order matters because Lean reads the file from top to bottom and accepts a name only after it has been defined. For example, the definition of \code{Inj} in Section~\ref{ch15:sec:injcomp} must come before \code{inj\_comp} and every later theorem that uses \code{Inj}. Following the order of the book guarantees this.

The examples are written for Lean version~4.19.0. Typing \code{lean -{}-version} in a terminal shows which version you have. Lean changes from one version to the next, so if an example that you copied exactly is rejected, compare your version with 4.19.0 first.

\subsection*{Running the check}
Open the file in VS Code. Lean checks it while you type, and when the cursor is inside a proof, the infoview beside the file shows the proof state at that point. To check the whole file at once from a terminal, move to the folder that contains it, as in Section~\ref{appD:sec:python}, and type
\par\noindent\begin{minipage}{\linewidth}
\begin{lstlisting}
lean proofs.lean
\end{lstlisting}
\end{minipage}\par
Lean reports each problem together with the line and column where it occurs. If the command prints nothing at all, every declaration was accepted, with no error and no warning. An \emph{error} means that the declaration containing it was rejected. A \emph{warning} means that the declaration was accepted but something in it deserves a second look. If you have copied every declaration, Lean prints two warnings, each followed by a note on how such warnings can be switched off. They say that the variables \code{hf} and \code{hg} in \code{surj\_comp\_assumed} are unused, and Section~\ref{ch16:sec:irrelevant} explains why that theorem never needs them.

\subsection*{Asking which axioms a proof uses}
To list the axioms on which a theorem depends (Section~\ref{ch16:sec:axioms}), add a line such as
\par\noindent\begin{minipage}{\linewidth}
\begin{lstlisting}
#print axioms MathCourse.surj_comp
\end{lstlisting}
\end{minipage}\par
after the line \code{end MathCourse}. The command uses the full name, with the prefix \code{MathCourse}, because it stands outside the namespace. Lean answers
\par\noindent\begin{minipage}{\linewidth}
\begin{lstlisting}
'MathCourse.surj_comp' does not depend on any axioms
\end{lstlisting}
\end{minipage}\par
and from now on the command \code{lean proofs.lean} prints this line as well. The theorems \code{inj\_comp} and \code{preimage\_meet} give the same answer. Not every theorem in the book does: \code{not\_surj\_succ} and \code{succ\_surj\_int} from Section~\ref{ch16:sec:debugging} depend on the standard axiom \code{propext}, through the library lemmas they use, and as that section explains, this is no cause for concern.

\subsection*{What an unfinished proof looks like}
To see how Lean reacts to a gap, temporarily replace the last line of the proof of \code{surj\_comp}, which is \code{rw [ha, hb]}, by the single word \code{sorry}, keeping the same indentation. Lean still accepts the file, but it prints the warning \code{declaration uses 'sorry'}, pointing at the name of the theorem, and the \code{\#print axioms} line above now reports
\par\noindent\begin{minipage}{\linewidth}
\begin{lstlisting}
'MathCourse.surj_comp' depends on axioms: [sorryAx]
\end{lstlisting}
\end{minipage}\par
Put the original line back before you rely on the file again: a file that contains \code{sorry} has not proved what it states.

\subsection*{Where to go next}
For further practice, \emph{Mathematics in Lean} is an interactive textbook with many exercises~\cite{mil}. It is built on Mathlib, the large mathematical library that the examples of this book do without, and its own instructions explain how to set it up. It is a natural next step once the statements, goals and proof states of Chapters~\ref{ch:lean1} and~\ref{ch:lean2} feel familiar.
